\documentclass[11pt]{article}
\usepackage{amsmath, amssymb, amsthm}

\newtheorem{proposition}{Proposition}
\newtheorem{corollary}{Corollary}
\newtheorem{lemma}{Lemma}
\theoremstyle{remark}
\newtheorem{remark}{Remark}

\newcommand{\Q}{\mathbb{Q}}

\title{Justification of the Symbolic Recurrence Method\\
as a Replacement for Lagrange Interpolation}
\author{}
\date{}

\begin{document}
\maketitle

\section*{Setup}

Fix rational numbers $j_1, j_2, j_3, \dots$ (in the application, these are
derived from the coefficients $\tau(p)$ of the discriminant cusp form, but
nothing below depends on their origin). Define a sequence $(h_n)_{n \ge 0}$
by the recurrence
\begin{equation}\label{eq:rec}
h_0 = 1, \qquad
h_n = \frac{1}{n}\left( j_n + \sum_{r=1}^{n-1} j_r\, h_{n-r} \right),
\qquad n \ge 1.
\end{equation}

The \emph{deformed} computation prepends a parameter to the $j$-sequence:
for $c_0 \in \Q$, apply \eqref{eq:rec} to the input sequence
$(c_0,\, j_1,\, j_2,\, \dots)$, i.e.\ with $j_1$ replaced by $c_0$ and the
original $j_r$ shifted to position $r+1$. Write $h_n(c_0)$ for the
resulting value, and define
\[
P_n(c_0) \;=\; n!\, h_n(c_0).
\]
The original numerical method recovers $P_n$ by evaluating $h_n(c_0)$ at
the integer nodes $c_0 = 1, 2, \dots, 2n+4$ and applying Lagrange
interpolation. The symbolic method instead runs \eqref{eq:rec} once over
the polynomial ring $\Q[c]$ with input sequence $(c,\, j_1,\, j_2,\, \dots)$.

\section*{The methods compute the same polynomial}

\begin{lemma}[Evaluation commutes with the recurrence]\label{lem:eval}
Let $H_n \in \Q[c]$ denote the output of recurrence \eqref{eq:rec} applied
over $\Q[c]$ to the input $(c,\, j_1,\, j_2,\, \dots)$. Then for every
$c_0 \in \Q$,
\[
H_n(c_0) = h_n(c_0).
\]
\end{lemma}

\begin{proof}
For fixed $c_0 \in \Q$, the evaluation map
$\mathrm{ev}_{c_0} \colon \Q[c] \to \Q$, $f \mapsto f(c_0)$,
is a ring homomorphism. The recurrence \eqref{eq:rec} is built exclusively
from ring operations: addition, multiplication, and multiplication by the
fixed scalars $1/n \in \Q$. A ring homomorphism (which here is moreover
$\Q$-linear) commutes with every expression formed from such operations.
Since $\mathrm{ev}_{c_0}$ sends the symbolic input sequence
$(c, j_1, j_2, \dots)$ to the numerical input sequence
$(c_0, j_1, j_2, \dots)$, induction on $n$ gives
$\mathrm{ev}_{c_0}(H_n) = h_n(c_0)$ for all $n$.
\end{proof}

Note that the argument uses no property of the data $j_1, j_2, \dots$
whatsoever; it is an identity in the inputs, valid for arbitrary rational
(indeed, arbitrary commutative-ring) values.

\begin{lemma}[Degree bound]\label{lem:deg}
$\deg_c H_n = n$, with leading term $c^n / n!$.
\end{lemma}

\begin{proof}
By induction. $H_0 = 1$ has degree $0$. In the deformed input, the
parameter $c$ occupies position $1$; every other input is a constant.
In \eqref{eq:rec}, the term $j_r h_{n-r}$ contributes $c$-degree at most
$1 + \deg H_{n-1}$ (when $r = 1$, since the deformed $j_1 = c$) and at most
$\deg H_{n-2}$ otherwise, while the lone term $j_n$ is constant in $c$ for
$n \ge 2$. Hence $\deg H_n \le 1 + \deg H_{n-1}$, and the unique
contribution attaining this bound is $\frac{1}{n}\, c\, H_{n-1}$. Its
leading coefficient is $\frac{1}{n} \cdot \frac{1}{(n-1)!} = \frac{1}{n!}$,
which is nonzero, so the degree is exactly $n$.
\end{proof}

\begin{proposition}[Equivalence]\label{prop:main}
For each $n \ge 2$, the polynomial $n!\, H_n \in \Q[c]$ produced by the
symbolic method equals the polynomial produced by Lagrange interpolation
of the values $P_n(c_0) = n!\, h_n(c_0)$ at the nodes
$c_0 = 1, 2, \dots, 2n+4$.
\end{proposition}

\begin{proof}
Let $L_n$ be the Lagrange interpolant: the unique polynomial of degree
$\le 2n+3$ satisfying $L_n(c_0) = P_n(c_0)$ at all $2n+4$ nodes. By
Lemma~\ref{lem:eval}, the polynomial $n!\, H_n$ also satisfies
$n!\, H_n(c_0) = n!\, h_n(c_0) = P_n(c_0)$ at every node, and by
Lemma~\ref{lem:deg} its degree is $n \le 2n+3$. Uniqueness of the
interpolant of degree $\le 2n+3$ through $2n+4$ points forces
$L_n = n!\, H_n$.
\end{proof}

\begin{corollary}
The two methods agree not only at the interpolation nodes but at every
$c_0 \in \Q$ (indeed, as elements of $\Q[c]$, coefficient by coefficient).
In particular, the symbolic method introduces no modeling assumption about
the data: it defers the evaluation step of the original computation rather
than altering it.
\end{corollary}

\begin{remark}
Proposition~\ref{prop:main} also explains an empirical observation of the
numerical method: the interpolants were always found to have degree exactly
$n$. By Lemma~\ref{lem:deg} this is a structural property of the recurrence
--- a consequence of the parameter entering linearly in position $1$ ---
and holds for arbitrary input data. The data determine the lower-order
coefficients of $H_n$, not its degree.
\end{remark}

\begin{remark}[Falsifiability]
The equivalence can be tested independently of the proof: evaluate
$n!\, H_n$ at points that were never interpolation nodes (e.g.\ negative
integers or non-integers such as $c_0 = -7$ or $3/2$) and compare with the
direct deformed computation $P_n(c_0)$; or compare the two polynomials
coefficient by coefficient with exact arithmetic. Agreement at non-node
points cannot be an artifact of fitting to the nodes.
\end{remark}

\end{document}
